\chapter{\texttt{tricycle/\_\_init\_\_.py}: registering the task}
\label{ch:task_init}

\textbf{Role:} gives the task its public name and tells Gymnasium and Isaac
Lab which three classes implement it. Runs once at start-up. 25 lines.

\codepart{\srcroot/luna_hifi_tasks/tricycle/__init__.py}{1}{11}

\begin{explain}
\lis{1}{11} Header comment. Lines 3--5 state the mechanism (entry point
  $\to$ import $\to$ registration); lines 7--8 are the two commands from the
  runbook, so anyone opening the file sees how to run it; lines 10--11 give
  the success criterion in the unit the reward is measured in (metres in
  2~s).
\end{explain}

\codepart{\srcroot/luna_hifi_tasks/tricycle/__init__.py}{12}{25}

\begin{explain}
\li{13} Gymnasium is the maintained fork of OpenAI Gym. Isaac Lab uses its
  registry as the phone book of tasks; nothing else from it is needed here.
\li{15} Imports the \code{agents} sub-package so that \code{agents.__name__}
  can be used on line 23 to spell the dotted path of the PPO config
  without hard-coding it. This also executes \srcpath{agents/__init__.py}, which
  is empty.
\li{17} \code{gym.register} adds one entry to the global registry. Everything
  inside the call is stored as data; no class is imported yet.
\li{18} The task id. Isaac Lab develop's own tasks drop the \code{-v0} suffix
  the old scripts used, and so do we. \code{--task Luna-Tricycle-Direct}
  on the command line must match this string exactly.
\li{19} The environment class as a dotted path string,
  \code{module:ClassName}. \code{__name__} is \srcpath{luna_hifi_tasks.tricycle},
  so this reads \srcpath{luna_hifi_tasks.tricycle.tricycle_env:TricycleEnv}.
  Gymnasium imports it lazily when \code{gym.make} is called
  (Section~\ref{sec:step_env}).
\li{20} Gymnasium normally wraps every environment in a checker that
  validates observation shapes against the declared space on each step. Isaac
  Lab environments return dictionaries of batched tensors, which the checker
  does not understand, so every Isaac Lab task disables it.
\li{21} Extra keyword arguments. Gymnasium passes them to the environment
  constructor, but Isaac Lab reads them earlier, from the registry, to find
  the config classes.
\li{22} \code{env_cfg_entry_point}: the environment config class. Isaac Lab's
  \code{load_cfg_from_registry(task, "env_cfg_entry_point")} imports
  \srcpath{luna_hifi_tasks.tricycle.tricycle_env_cfg} and instantiates
  \code{TricycleEnvCfg} (Section~\ref{sec:step_cfg}).
\li{23} \srcpath{rsl_rl_cfg_entry_point}: the agent config for the
  \code{rsl_rl} library, built from \code{agents.__name__}
  (\srcpath{luna_hifi_tasks.tricycle.agents}) plus the module and class. The key
  name is a convention: \code{--rl_library rsl_rl} makes Isaac Lab look up
  \srcpath{rsl_rl_cfg_entry_point}; a \code{skrl} backend would look for
  \srcpath{skrl_cfg_entry_point}, and so on.
\li{24} Closes the \code{kwargs} dictionary.
\li{25} Closes the \code{gym.register} call.
\end{explain}

\begin{isaacbox}{the three classes behind one task id}
Every Isaac Lab task is a triple: an \emph{env class} (behaviour), an
\emph{env cfg class} (data: scene, physics, spaces), and one \emph{agent cfg
class per learning library}. The registry stores their import paths as
strings so that listing tasks does not import heavy modules. The
\code{gym.make} call and the two \srcpath{load_cfg_from_registry} calls are where
the strings are turned into classes.
\end{isaacbox}
